\documentclass[11pt,a4paper]{article}
\usepackage[T1]{fontenc}\usepackage[utf8]{inputenc}\usepackage{lmodern}
\usepackage[a4paper,margin=2.2cm]{geometry}
\usepackage{amsmath,amssymb}\usepackage{enumitem}\usepackage{booktabs,longtable,array}
\usepackage[hidelinks]{hyperref}
\setlength{\parindent}{0pt}\setlength{\parskip}{0.4em}
\newtheorem{sublemma}{Sublemma}
\title{\textbf{OP3 / B2: merged versus split breakdown surfaces ---
a conditional sublemma}\\\large Recon v1 --- proposal-only}
\author{Per reviewer directive after acceptance of F7 rev.~1.1 and the BB
bridge}
\date{12 June 2026 --- rev.~1.1 (review safeguards + H2 verification
outcome applied)}
\begin{document}\maketitle

\section*{0. Scope and firewall}
B2 formalises Bridge~2 of the BB-bridge recon as a conditional sublemma at
the level of the recorded breakdown criterion. Everything is inherited-Level-B
conditional (kernel route); the statement concerns the \emph{structure of the
ordered-branch validity boundary}, not the transition dynamics, and respects
the recorded interface requirements. Proposal-only; no preprint edits.

\section{Setting and hypotheses}
Recorded breakdown indicators of the ordered geometric branch:
(i)~amplitude collapse $\rho\to0$;
(ii)~$T_{\rm loc}\to T_c$;
(iii)~curvature/gradient invariants reaching the condensate threshold
$\ell_{\rm cond}\sim m_\sigma^{-1}$, i.e.\ a curvature scale
$K\sim m_\sigma^{2}$.
The record already ties (iii) to the UV threshold: the same radial mode
$m_\sigma$ sets both. Hypotheses:
\begin{description}[leftmargin=2.0em,itemsep=2pt]
\item[(H1)] the Level-B induced-stiffness kernel holds, with
$\gamma_{\hat a_{\rm eff}}=0$ and no additional gradient-sector scale (the
F7 assumptions), so the $\sigma$-gap follows one of the two F7 branches:
gradient-native $m_\sigma^2=c\,u$, or potential-imported pinned $m_\sigma^2$;
\item[(H2)] the recorded indicator~(i), $\rho\to0$, and the
gradient-amplitude collapse $u\to0$ are treated as \emph{parallel
amplitude-loss indicators} of the same ordered-regime failure, not
necessarily the same variable. (Verification audit, ledger of this date:
the standalone amplitude symbol $\rho$ is recorded primarily for the
coherent-sector order parameter $\Phi_{\rm eff}=\rho e^{i\theta}$ and is
not defined at the breakdown-criterion site; the sublemma below therefore
phrases its conclusions against $u\to0$ directly and uses indicator~(i)
only as the recorded parallel signal.)
\item[(H3)] trajectories are characterised by their curvature/gradient
invariant $K$, bounded below by some $K_0>0$ on the approach considered.
\end{description}

\section{Sublemma and proof sketch}
\begin{sublemma}[conditional; merged vs split breakdown surfaces]
Under (H1)--(H3):
(a) On the gradient-native branch, indicator (iii) fires at
$u_*\simeq K_0/c>0$, strictly before amplitude collapse; the locus
$u\to0$ lies outside the ordered-EFT validity domain for trajectories
satisfying (H3), and the surfaces (i) and (iii) merge into a single breakdown
boundary in the slowly-varying limit.
(b) On the potential-imported branch, the (iii)-threshold is
$u$-independent; every trajectory with $K<m_\sigma^2$ reaches arbitrarily
small $u$ without triggering indicator (iii), so the surfaces (i) and (iii)
split, and an intermediate regime exists (small amplitude, indicator (iii)
not yet triggered; other recorded indicators may still apply).
\end{sublemma}
\emph{Proof sketch.} Indicator (iii) is the condition $K\gtrsim
m_\sigma^2(u)$. (a)~Gradient-native: $m_\sigma^2(u)=c\,u$ decreases with the
order parameter, so the threshold curvature falls linearly with $u$; for
$K\ge K_0$ the condition is met at $u_*\simeq K_0/c>0$. The cutoff thus
``chases'' the order parameter down: $\Lambda_{\rm UV}(u)=m_\sigma(u)=
\sqrt{c\,u}\to0$, and no EFT-internal account of the collapse point exists.
(b)~Potential-imported: $m_\sigma^2$ is fixed; the condition
$K\gtrsim m_\sigma^2$ is $u$-independent, so low-curvature trajectories
($K<m_\sigma^2$) do not trigger indicator (iii) down to
$u\to0$.\hfill$\square$

\section{Corollary: a transition-side discriminator, and one sharp question}
\textbf{Discriminator.} The two F7 branches predict structurally different
approaches to the pre-ordered regime. Merged (gradient-native): the
transition region is reached only through total breakdown of the ordered
EFT --- there is \emph{no semiclassical window} in which the ordered-branch
description covers the small-amplitude regime. Split (potential-imported):
a window exists in which $u\approx0$ physics does not yet trigger
indicator (iii) below the fixed cutoff (other recorded indicators may still
apply).
\textbf{B2-Q1 (flagged question, not a claim).} The recorded
critical-fluctuation route to the primordial tilt ($n_s\sim0.96$--$0.97$
from critical $O(4)$ fluctuation estimates) computes near the transition.
Whether that route presupposes the split-branch semiclassical window, or is
formulated at the pre-emergent Euclidean level where the present sublemma
does not apply, is a sharp follow-up question that would couple the F7
branch choice to a recorded observable estimate. A first reading of the
recorded construction (\S sec:primordial\_estimates, this date) localises
the question: the critical correlator and its log-corrected $G(r)$ live at
the pre-emergent Euclidean level, where the present sublemma is silent,
while the conversion step $k^2\Phi_k\sim G_N\delta\rho_k$ and the
freeze-out discipline use post-onset structure --- B2-Q1 therefore
concentrates on the conversion/freeze-out step, not on the correlator
itself. The full answer still requires the dedicated audit, not assumption.
\textbf{Late-time invisibility.} Both branches are identical at $u\simeq
u_0$ (heavy $m_\sigma$, far above Hubble scales); the discriminator lives
only near the threshold. No observable consequence at current sensitivity
follows --- consistent with the F7 classification discipline.

\section{Status and verdict}
\begin{enumerate}[leftmargin=1.9em,itemsep=2pt,label=\textbf{V\arabic*.}]
\item Bridge~2 is now a conditional sublemma with explicit hypotheses
(H1)--(H3) and an elementary proof sketch; inherited level: Level~B
conditional (kernel route), candidate status pending review.
\item Promotion path: (a)~promotion of the kernel lemma itself;
(b)~the H2 verification audit is \emph{performed} (ledger of this date):
the standalone $\rho$ is recorded for the coherent-sector order parameter
and is undefined at the breakdown-criterion site, so the parallel-indicator
formulation of (H2) replaces the variable identification; the remaining
promotion need is (a) alone.
\item The sublemma equips OP3 with a structural dichotomy whose resolution
simultaneously settles F7, and isolates B2-Q1 as the one place where the
dichotomy may touch a recorded observable estimate. No preprint edits
before review.
\end{enumerate}
\end{document}
